\chapter{Introduction}\label{ch:intro}

\section{Problem and motivation}
An applicant for a Schengen short-stay visa in Türkiye prepares a file of about a dozen documents:
a passport, a photograph, an application form, travel and accommodation bookings, insurance, bank
statements and proof of employment or study. What exactly is required is spread over several
official pages. A Commission decision lists the supporting documents common to all Schengen
states for applicants in Türkiye~\autocite{eu2021harmonised}, and each state adds its own
checklist with its own wording, limits and exceptions. Table~\ref{tab:nl-no} shows how two states
differ on the same points: Norway asks for translations into English or Norwegian and refuses
supporting documents older than six months, while the checklist of the Netherlands states neither
rule~\autocite{nl2025checklist,no2025checklist}.

The documents must also agree with each other and with the trip. The insurance has to cover every
day of the stay, the leave in the employer's letter has to cover the trip, the passport has to
remain valid for three months after the return, and the name on a booking has to be the name in
the passport. None of these checks is difficult, but there are many of them, they depend on dates
that change when the trip changes, and a single miss can mean a refused or delayed application
after the fee has been paid and an appointment has been used. Applicants who do not pay an agency
make these checks by hand.

A general-purpose AI assistant does not solve this. It answers fluently, but it does not show
which official sentence a requirement comes from, it cannot tell whether that sentence is still in
force, and it has no view of the applicant's whole file. VisaAssist AI takes the opposite route:
requirements are written down once, from official sources, with the passage they rely on; the
applicant's documents are checked against them by rules that give the same answer every time; and
every result points back to its source. The system helps to prepare a file. It does not predict
or influence the decision, which rests with the consulate.

\section{Objectives}
The project pursues the five objectives of the approved proposal:
\begin{description}
  \item[H1 Personalisation] a checklist for the applicant's destination, visa type, profile and dates;
  \item[H2 Document pre-check] missing, outdated and mutually inconsistent documents are found before submission;
  \item[H3 Up-to-date information] changes of official sources are detected and reviewed;
  \item[H4 Source-grounded information] every requirement, finding and answer cites an official source;
  \item[H5 Measurable AI performance] the AI components are evaluated against baselines.
\end{description}

\section{Research questions}\label{sec:rq}
\begin{description}
  \item[AS1] Does a hybrid approach (vision-language model and rule engine) detect missing and
    inconsistent documents in multi-document application files better than an LLM-only approach,
    per error type (precision, recall, F1)?
  \item[AS2] How does the hybrid system compare with human reviewers in accuracy and time?
  \item[AS3] How accurately does source monitoring separate meaningful requirement changes from
    cosmetic page changes?
\end{description}

\section{Scope}
% ADR 0001 and its amendment.
The pilot covers the Schengen short-stay visa (type C) for tourism, applied for in Türkiye, for
employed and student applicants. The Netherlands is the main destination; Norway is added as a
second national layer because its checklist states a translation requirement. Predicting visa
decisions, legal advice and submitting applications on the applicant's behalf are out of scope.

\section{Contributions}
Chapter~\ref{ch:background} shows that the building blocks exist separately: document
understanding models extract fields, rule frameworks validate identity documents, and
retrieval-augmented systems cite passages. The project does not claim any of these as new. Its
contribution is their combination for one concrete task, and the measurement of that combination:
\begin{itemize}
  \item \textbf{Versioned, source-linked requirements.} Every checklist item and every rule
    stores the official source, the section and the passage it relies on, and belongs to a
    requirement set identified by the hash of its files, so an earlier check stays explainable
    after the rules change.
  \item \textbf{Checks across a multi-document file.} Missing documents and disagreements between
    documents (names, dates, amounts, periods) are found on the fields of visual documents such as
    scans and photographs, not on text that is already machine-readable.
  \item \textbf{Change detection that reaches only the affected applicants.} Official pages are
    snapshotted weekly; a change approved by a reviewer becomes a new requirement set, and only
    the files whose results change are notified.
  \item \textbf{A measured comparison.} The hybrid system is compared with an LLM-only baseline
    and with human reviewers on a synthetic data set with controlled errors
    (Section~\ref{sec:rq}).
\end{itemize}
The fall semester delivers the first two items without the AI layer: the applicant enters the
fields, and the rule engine checks them. Field extraction, source monitoring with review, and the
evaluation follow in the spring semester.

\section{Organisation of the report}
Chapter~\ref{ch:background} describes the visa process and the related work.
Chapter~\ref{ch:requirements} states the functional and non-functional requirements.
Chapter~\ref{ch:design} presents the architecture, the data model and the rule kinds.
Chapter~\ref{ch:implementation} describes what has been built so far.
Chapter~\ref{ch:evaluation} sets out how the system is tested and how the research questions will
be answered. Chapter~\ref{ch:management} covers the plan, the risks and the progress against the
plan. Chapter~\ref{ch:ethics} discusses personal data, consent and the limits of what the system
may claim. Chapter~\ref{ch:conclusion} summarises the state of the work and the plan for the
spring semester.
